\documentclass[../../main.tex]{subfiles}
\begin{document}

\begin{center}
  {\fontsize{48pt}{58pt}\selectfont \textbf{2004}\normalsize\textbf{年度}}
\end{center}

\begin{table}[htb]
  \centering
  \caption{出題テーマと難易度}
  \begin{tabular}{|c|p{9cm}|c|}
    \hline
    \textbf{問題番号} & \textbf{テーマ} & \textbf{難易度} \\
    \hline
    第1問 & 区分的な関数を含む定積分の極限 & \\
    \hline
    第2問 & 複素平面上の帯状の領域 & \\
    \hline
    第3問 & 加法性を持つ写像と有理数倍 & \\
    \hline
    第4問 & 仰角から求める棒の高さ & \\
    \hline
    第5問 & 不等式を満たす自然数の組の数 & \\
    \hline
    第6問 & 円内の格子正方形の漸近数 & \\
    \hline
  \end{tabular}
\end{table}

\begin{table}[htb]
  \centering
  \caption{解答}
  \renewcommand{\arraystretch}{1.4}
  \begin{tabular}{|c|c|p{9cm}|}
    \hline
    \textbf{問題番号} & \textbf{小問} & \textbf{答え} \\
    \hline
    第1問 & -- & $\dfrac{1}{4}$ \\
    \hline
    第2問 & -- & $\alpha=X+Yi$に対し$|X-Y+1|\le\sqrt{2}$（境界を含む） \\
    \hline
    第3問 & -- &  \\
    \hline
    第4問 & -- & 約$28\,\mathrm{m}$ \\
    \hline
    第5問 & -- & $\dfrac{n^2(n-1)(n+1)}{12}$通り \\
    \hline
    第6問 & -- & $\displaystyle\lim_{n\to\infty}\dfrac{N(n)}{n^2}=\pi$（証明問題） \\
    \hline
  \end{tabular}
\end{table}

\paragraph{試験形態}
6問．試験時間は要確認．

\end{document}
